\documentclass[11pt,
			technote,
			%draft
			]{article}
%\renewcommand{\rmdefault}{ptm} %change la police par d�faut
\usepackage{geometry}                % See geometry.pdf to learn the layout options. There are lots.
\usepackage[applemac]{inputenc} 

\usepackage{tikz}  %for plotting cute graphs
\usetikzlibrary{plotmarks}
\usepackage{tikz} \usepackage{pgfplots}


\usepackage{subfig}% for table side by side

\usepackage{amsmath}
\usepackage[T1]{fontenc}  
\usepackage{graphicx}
\usepackage{lscape}
\usepackage{appendix} %for using appendices

\usepackage{multirow} % for using multirows in table
%\usepackage[table]{xcolor}  %for using alternate column colors for each row in a table

\usepackage{listings}	%for using code listings
%\geometry{letterpaper}                   % ... or a4paper or a5paper or ... 
%\geometry{landscape}                % Activate for for rotated page geometry
%\usepackage[parfill]{parskip}    % Activate to begin paragraphs with an empty line rather than an indent
\usepackage{graphicx}
\usepackage{amssymb}
\usepackage{epstopdf}

\usepackage{array,booktabs} %for arrays
\usepackage{longtable} %for long table on multiple page

\DeclareGraphicsRule{.tif}{png}{.png}{`convert #1 `dirname #1`/`basename #1 .tif`.png}
                       
                       
                       
 %\includeonly{
%03-introduction/introduction,
%04-related_work/related_work,
%05-folding_operations_specifications/folding_operations_specifications,
%06-prototyping_tool_specifications/prototyping_tool_specifications,
%07-rendering_prototypes/rendering_prototypes,
%08-conclusion/conclusion,
%}

                       
\begin{document}
%%%%%%%%%%%%%%%%%%%%%%%%%%%%%%%%%%%%%%%%%
% University Assignment Title Page 
% LaTeX Template
%
% This template has been downloaded from:
% http://www.latextemplates.com
%
% Original author:
% WikiBooks (http://en.wikibooks.org/wiki/LaTeX/Title_Creation)
% 
% Instructions for using this template:
% This title page is presently capable of being compiled as is. This is not 
% useful for including it in another document. To do this, you have two options: 
%
% 1) Copy/paste everything between \begin{document} and \end{document} 
% starting at \begin{titlepage} and paste this into another LaTeX file where you 
% want your title page.
% OR
% 2) Remove everything outside the \begin{titlepage} and \end{titlepage} and 
% move this file to the same directory as the LaTeX file you wish to add it to. 
% Then add \input{./title_page_1.tex} to your LaTeX file where you want your
% title page.
%
%%%%%%%%%%%%%%%%%%%%%%%%%%%%%%%%%%%%%%%%%

%----------------------------------------------------------------------------------------
%	PACKAGES AND OTHER DOCUMENT CONFIGURATIONS
%----------------------------------------------------------------------------------------


\begin{titlepage}

\newcommand{\HRule}{\rule{\linewidth}{0.5mm}} % Defines a new command for the horizontal lines, change thickness here

\center % Center everything on the page
 
%----------------------------------------------------------------------------------------
%	HEADING SECTIONS
%----------------------------------------------------------------------------------------

%\textsc{\LARGE Universit� Catholique de Louvain}\\[1.5cm] % Name of your university/college
%\textsc{\Large Ecole Polytechnique de Louvain}\\[0.5cm] % Major heading such as course name

\begin{minipage}{6in}
  \centering
  $\vcenter{\hbox{\includegraphics[height=100px]{cover/img/ucl-logo}}}$
  %\hspace*{1cm}
  \hfill
  $\vcenter{\hbox{\includegraphics[height=60px]{cover/img/epl-logo}}}$
\end{minipage}


\vspace{4cm}
\textsc{\Large LSINF2335: Programing paradigms: theory, practice and applications}\\[0.5cm] % Major heading such as course name



%----------------------------------------------------------------------------------------
%	TITLE SECTION
%----------------------------------------------------------------------------------------

\HRule \\[0.4cm]
{ \huge \bfseries Usage contracts for Ruby}\\[0.4cm] % Title of your document
\HRule \\[1.5cm]
 
%----------------------------------------------------------------------------------------
%	AUTHOR SECTION
%----------------------------------------------------------------------------------------

\begin{minipage}{0.4\textwidth}
\begin{flushleft} \large
\emph{Author:}\\
No� \textsc{Phan}\\ % Your name
Xavier \textsc{Crochet}
\end{flushleft}
\end{minipage}
~
\begin{minipage}{0.4\textwidth}
\begin{flushright} \large
\emph{Teachers:} \\
Kim \textsc{Mens}\\ % Supervisor's Name
Sebastian \textsc{Andres Gonzalez Montesinos} % Supervisor's Name
\end{flushright}
\end{minipage}\\[4cm]

% If you don't want a supervisor, uncomment the two lines below and remove the section above
%\Large \emph{Author:}\\
%John \textsc{Smith}\\[3cm] % Your name

%----------------------------------------------------------------------------------------
%	DATE SECTION
%----------------------------------------------------------------------------------------

{\large \today}\\[3cm] % Date, change the \today to a set date if you want to be precise

%----------------------------------------------------------------------------------------
%	LOGO SECTION
%----------------------------------------------------------------------------------------
\vfill % Fill the rest of the page with whitespace


 
%----------------------------------------------------------------------------------------



\end{titlepage}

%\maketitle

\renewcommand{\thepage}{\Roman{page}} %change the numbering style


\tableofcontents
\clearpage
\setcounter{page}{1}
\renewcommand{\thepage}{\arabic{page}} %change the numbering style
\newpage



\section{Introduction}
In this report we will explain the structure of the application designed to support the usage contracts for Ruby. First we present the overview of the applications and the concepts used to implements the usage contracts. The then structure of the DSL is detailed. Finally, a discussion about limitations and improvements that could be done is made.
\section{overview}
content
\section{DSL specification}
content
\section{Improvements}
\begin{itemize}
\item{Give more informations about the contracts failures such as a line number where the contract is not respected or about the condition not met.}
\item{Let users target a type of methods in contracts, such as accessors, mutator or initializations methods}
\item{Let users specify several methods within one contract}
\end{itemize}  
 
\section{Conclusion}
content




\end{document}